\documentclass{article}
\usepackage[UTF8]{ctex}
\usepackage{amsmath,amssymb,amsfonts}
\usepackage{enumitem}
\usepackage{geometry}
\geometry{a4paper, margin=1in}
% === 错题本排版优化（不参与题目标识） ===
\setlength{\parindent}{0pt}
\setlength{\parskip}{0.35em plus 0.10em minus 0.08em}
\setlength{\emergencystretch}{3em}
\setlist{itemsep=0.25em, topsep=0.40em, parsep=0pt}
\allowdisplaybreaks[2]
% === 错题本排版优化结束 ===
\begin{document}
\section*{中南大学2026年数学分析真题}
\section*{146．中南大学2026年数学分析试卷（回忆版）}

\subsection{题目1}
1．设 \(x=1+t^{3}, y=\ln (1+t)\) ，求 \(\left.\frac{\mathrm{d} y}{\mathrm{~d} x}\right|_{t=1}\) ．

\subsection{题目2}
2．求极限 \(\lim _{x \rightarrow 0} \frac{\int_{x}^{x^{3}} \ln \left(1+t^{2}\right) \mathrm{d} t}{x^{3}}\) ．

\subsection{题目3}
3 ．求积分 \(\int_{0}^{+\infty} \frac{2 x \ln \left(1+x^{2}\right)}{\left(1+x^{2}\right)^{2}} \mathrm{~d} x\) ．

\subsection{题目4}
4．求级数 \(\sum_{n=1}^{\infty} \frac{1}{n(n+2)}\) 的和。

二．（20 分）设 \(f:[0,1]^{3} \rightarrow \mathbb{R}\) 为连续函数，证明：\(f\) 的像为有界闭区间．

三．（20 分）设 \(f(x)\) 在［0，1］上连续，在 \((0,1)\) 上二阶可导，\(f(0)=0, f\left(\frac{1}{2}\right)=\frac{1}{2}, f(1)=0\) ，证明：存在\(\xi \in(0,1)\) ，满足 \(f^{\prime \prime}(\xi)=-4\) ．

四．（20 分）求积分

\[
\iint_{\sum} x \mathrm{~d} y \mathrm{~d} z+y \mathrm{~d} z \mathrm{~d} x+z \mathrm{~d} x \mathrm{~d} y .
\]

其中 \(\sum\) 为球面 \(x^{2}+y^{2}+z^{2}=25\) ，取曲面的内侧．

五。（20分）设 \(f(x)\) 在 \([1,+\infty)\) 上一致连续，\(\varphi(x)\) 在 \([1,+\infty)\) 上连续，且 \(\lim _{x \rightarrow+\infty}[f(x)-\varphi(x)]=0\) ，证明：\(\varphi(x)\) 在 \([1,+\infty)\) 上一致连续。

六．（10 分）求反常积分 \(\int_{0}^{+\infty} \frac{e^{-a x} \sin b x}{x} \mathrm{~d} x ; a>0, b \in \mathbb{R}\) ．

\begin{itemize}
\end{itemize}

\subsection{题目10}
原题边界“t．”：（10 分）设 \(S(a)=\sum_{n=1}^{\infty} \frac{\sin ^{2}(n \pi a)}{n^{3}}, a \in[0,1]\) ．

\begin{itemize}
\item[（1）] 求 \(S(a)\) 的收敛域，并证明在收敛域内 \(S(a)\) 一致收敛。
  \item[（2）] 证明：\(S(a)\) 在 \([0,1]\) 上连续可微．
  \item[（3）] 证明：\(f(a)=\frac{S(a)}{a^{2}}\) 单调递减．提示：\(S^{\prime \prime}(a)=\sum_{n=1}^{\infty} \frac{2 \pi^{2} \cos (2 n \pi a)}{n}=-\pi^{2} \ln \left(4 \sin ^{2}(\pi a)\right)\) ．
\end{itemize}

\subsection{题目11}
原题拆分片段：
几（10 分）已知在 \(\mathbb{R}^{3}\) 中，在点 \(\mathbf{x}_{0}\) 处存在某个电子，由该电子产生的电场为 \(\mathbf{E}=\frac{q}{4 \pi \varepsilon_{0}} \frac{\mathbf{x}-\mathbf{x}_{0}}{\left\|\mathbf{x}-\mathbf{x}_{0}\right\|^{3}}\) ，其中\(x \in \mathbb{R}^{3}\) ，且 \(x \neq x_{0}, \varepsilon_{0}\) 是正常数．设 \(\Omega\) 为 \(\mathbb{R}^{3}\) 中的单连通区域，其边界 \(\partial \Omega\) 为光滑曲面．

\begin{itemize}
  \item[（1）] 求证：\(x_{0} \notin \Omega\) 时，电通量 \(\Phi=\iint_{\partial \Omega} \mathrm{Ed} S=0\) ．
  \item[（2）] \(x_{0} \in \Omega\) 时，电通量 \(\Phi=\iint_{\partial \Omega} \mathrm{Ed} S\) 为常数，求此常数．
\end{itemize}
\end{document}
